% 03-architecture.tex
\chapter{System Architecture}
\label{ch:architecture}

\section{Overview}
\label{sec:arch-overview}

Doki's architecture follows a client-server model with three primary
layers: the CLI client (\texttt{doki}), the daemon (\texttt{dokid}),
and the container runtime. The CLI communicates with the daemon over
a Unix domain socket (or TCP on Android), and the daemon manages
container lifecycle operations through a pluggable runtime system.

\begin{figure}[htbp]
  \centering
  \begin{tikzpicture}[
    node distance=0.5cm,
    every node/.style={font=\sffamily\footnotesize},
  ]
    % User layer
    \node[comp, minimum width=3cm] (cli) {doki CLI};

    % API layer
    \node[compAlt, below=0.6cm of cli] (api) {REST API\\{\scriptsize\color{textgray} Unix Socket / TCP}};

    % Daemon layer
    \node[comp, below=0.6cm of api] (daemon) {dokid Daemon\\{\scriptsize\color{textgray} State Management + DNS}};

    % Runtime layer
    \node[compTeal, below left=0.6cm and 0.3cm of daemon] (proot) {proot\\{\scriptsize\color{textgray} ptrace}};
    \node[compTeal, below right=0.6cm and 0.3cm of daemon] (ns) {Namespaces\\{\scriptsize\color{textgray} unshare}};

    % Kernel
    \node[compPurple, below=0.6cm of daemon] (kernel) {Linux Kernel\\{\scriptsize\color{textgray} syscalls + cgroups}};

    % Arrows
    \draw[arr] (cli) -- (api);
    \draw[arr] (api) -- (daemon);
    \draw[arrThick] (daemon) -| (proot);
    \draw[arrThick] (daemon) -| (ns);
    \draw[arrDash] (proot) -- (kernel);
    \draw[arrDash] (ns) -- (kernel);

    % Regions
    \node[region, fit=(cli)(api), label={[annotBold]above:User Space}] {};
    \node[region, fit=(daemon), label={[annotBold]right:Daemon}] {};
    \node[region, fit=(proot)(ns), label={[annotBold]below:Runtime Selection}] {};
  \end{tikzpicture}
  \caption{Doki system architecture. The daemon selects the appropriate
    runtime based on available host capabilities.}
  \label{fig:architecture}
\end{figure}

\section{Daemon Architecture}
\label{sec:daemon}

The \texttt{dokid} daemon is the central component of the Doki system.
It runs as a background process and manages all container lifecycle
operations. The daemon is implemented in Go and uses the standard
\texttt{net/http} package for its REST API.

\subsection{Startup Sequence}
\label{sec:startup-sequence}

When \texttt{dokid} starts, it performs the following initialization
steps:

\begin{enumerate}[leftmargin=*, itemsep=3pt]
  \item Load configuration from \texttt{\$DOKI\_DATA\_DIR/config.json}
  \item Initialize the container state store (JSON-based persistence)
  \item Set up the internal DNS server on \texttt{127.0.0.11:8053}
        (Android) or \texttt{127.0.0.11:53} (Linux)
  \item Probe host capabilities and select strongest isolation runtime
  \item Initialize DokiLink-Lite mesh (Ed25519 keypair + CA generation)
  \item Register container state with the DNS server
  \item Start the REST API listener
\end{enumerate}

\subsection{State Management}
\label{sec:state-management}

Container state is persisted as JSON files in the
\texttt{\$DOKI\_DATA\_DIR/containers/} directory. Each container has
its own state file containing:

\begin{itemize}[leftmargin=*, itemsep=2pt]
  \item Container ID (SHA-256 hash of configuration)
  \item Name, image, and command
  \item Current state (created, running, stopped, exited)
  \item Network configuration (IP address, veth pair names)
  \item PID of the container process
  \item Exit code and timestamps
  \item Mount points and volume mappings
\end{itemize}

The daemon uses a \texttt{signal(0)} polling mechanism to detect
process termination, replacing the earlier \texttt{ExitChan} approach
that suffered from channel disconnect races.

\section{CLI Architecture}
\label{sec:cli}

The \texttt{doki} CLI is a single binary that provides 244 commands
across 9 categories. It communicates with the daemon via the REST API
and handles user input parsing, output formatting, and error reporting.

\subsection{Command Structure}
\label{sec:command-structure}

Doki's CLI follows a hierarchical command structure:

\begin{codebox}[title={CLI command hierarchy}]{text}
doki <command> [subcommand] [flags] [arguments]

doki container create   --name web --image alpine:3.18
doki container ls       --format "table {{.ID}}\t{{.Names}}"
doki image pull         nginx:latest
doki compose up         -d --build
doki mesh status
doki link add           mypeer 192.168.1.42:7432
\end{codebox}

\subsection{Flag Parsing}
\label{sec:flag-parsing}

Doki uses a two-pass flag parsing system. Global flags (such as
\texttt{-d}, \texttt{-q}, \texttt{-f}, \texttt{-p}, and
\texttt{--profile}) are recognized both before and after the
subcommand. This allows natural command ordering:

\begin{codebox}[title={Equivalent flag positions}]{text}
doki compose -f docker-compose.yml -d up
doki compose up -f docker-compose.yml -d
\end{codebox}

The \texttt{flagsWithValue} map tracks flags that consume the next
argument as their value. Ambiguous single-letter flags (\texttt{-f},
\texttt{-t}, \texttt{-s}) are excluded from this map to prevent
conflicts with container IDs and other positional arguments.

\section{Communication Protocol}
\label{sec:protocol}

The CLI and daemon communicate over a REST API using HTTP/1.1. The
API follows RESTful conventions with standard HTTP methods and status
codes.

\begin{table}[htbp]
  \centering
  \caption{API endpoint examples and their HTTP methods.}
  \label{tab:api-endpoints}
  \begin{tabularx}{\textwidth}{l l X}
    \toprule
    \textbf{Method} & \textbf{Endpoint} & \textbf{Description} \\
    \midrule
    GET    & \texttt{/containers/json}       & List all containers \\
    POST   & \texttt{/containers/create}     & Create a new container \\
    GET    & \texttt{/containers/\{id\}/json} & Inspect a container \\
    POST   & \texttt{/containers/\{id\}/start} & Start a container \\
    POST   & \texttt{/containers/\{id\}/stop}  & Stop a container \\
    DELETE & \texttt{/containers/\{id\}}     & Remove a container \\
    POST   & \texttt{/containers/\{id\}/exec} & Execute command in container \\
    GET    & \texttt{/images/json}           & List all images \\
    POST   & \texttt{/images/create}         & Pull or build an image \\
    GET    & \texttt{/\_ping}                & Health check \\
    \bottomrule
  \end{tabularx}
\end{table}

\section{Container Lifecycle}
\label{sec:lifecycle}

A container in Doki follows a well-defined state machine with four
primary states: \texttt{created}, \texttt{running},
\texttt{stopped}, and \texttt{exited}. The transitions between
states are managed by the daemon and triggered by CLI commands.

\begin{figure}[htbp]
  \centering
  \begin{tikzpicture}[
    node distance=2.2cm,
    on grid, auto,
    every state/.style={
      rectangle, rounded corners=4pt,
      draw=nodeblue!70, fill=fillblue,
      minimum width=2.2cm, minimum height=0.9cm,
      font=\sffamily\footnotesize, text=textdark,
      align=center, inner sep=5pt
    },
    initial/.style={
      initial text={},
      initial distance=1.5cm
    },
    accepting/.style={
      accepting by double,
      double distance=2pt,
      inner sep=1pt
    },
    every edge/.style={
      draw=textdark!70,
      line width=0.5pt,
      -{Stealth[length=2.5mm, width=2mm]},
      font=\sffamily\tiny,
      node distance=1cm
    }
  ]
    \node[state, initial] (created) {created};
    \node[state, right=of created] (running) {running};
    \node[state, below=of running] (stopped) {stopped};
    \node[state, accepting, left=of stopped] (exited) {exited};

    \path[->]
      (created) edge node {start} (running)
      (running) edge node {stop} (stopped)
      (running) edge[bend left=20] node {kill} (exited)
      (stopped) edge node {start} (running)
      (stopped) edge node {rm} (exited);
  \end{tikzpicture}
  \caption{Container lifecycle state machine. States are shown as
    rectangles; transitions are labeled with the triggering command.}
  \label{fig:lifecycle}
\end{figure}

\section{Isolation Runtime Selection}
\label{sec:runtime-selection}

Doki supports 12 isolation levels, implemented as separate runtime
modules in \texttt{pkg/runtime/runners/}. At startup, the daemon's
\texttt{RuntimeRegistry} probes the host system to determine which
runtimes are available and selects the strongest one.

\begin{figure}[htbp]
  \centering
  \begin{tikzpicture}[
    node distance=0.35cm,
    every node/.style={font=\sffamily\footnotesize},
  ]
    % Isolation levels (bottom = strongest)
    \node[rectangle, rounded corners=2pt, draw=nodepurple!60, fill=nodepurple!10, minimum width=6cm, minimum height=0.55cm, text=textdark, align=center, inner sep=3pt] (L1) at (0, 0.65) {microVM};
    \node[rectangle, rounded corners=2pt, draw=nodelteal!60, fill=nodelteal!10, minimum width=6cm, minimum height=0.55cm, text=textdark, align=center, inner sep=3pt] (L2) at (0, 1.30) {gVisor};
    \node[rectangle, rounded corners=2pt, draw=nodeblue!60, fill=nodeblue!10, minimum width=6cm, minimum height=0.55cm, text=textdark, align=center, inner sep=3pt] (L3) at (0, 1.95) {Namespaces};
    \node[rectangle, rounded corners=2pt, draw=nodeorange!60, fill=nodeorange!10, minimum width=6cm, minimum height=0.55cm, text=textdark, align=center, inner sep=3pt] (L4) at (0, 2.60) {Chroot};
    \node[rectangle, rounded corners=2pt, draw=nodeblue!60, fill=nodeblue!10, minimum width=6cm, minimum height=0.55cm, text=textdark, align=center, inner sep=3pt] (L5) at (0, 3.25) {Proot};
    \node[rectangle, rounded corners=2pt, draw=nodelteal!60, fill=nodelteal!10, minimum width=6cm, minimum height=0.55cm, text=textdark, align=center, inner sep=3pt] (L6) at (0, 3.90) {PKdroid};
    \node[rectangle, rounded corners=2pt, draw=nodeorange!60, fill=nodeorange!10, minimum width=6cm, minimum height=0.55cm, text=textdark, align=center, inner sep=3pt] (L7) at (0, 4.55) {QEMU-user};
    \node[rectangle, rounded corners=2pt, draw=nodepurple!60, fill=nodepurple!10, minimum width=6cm, minimum height=0.55cm, text=textdark, align=center, inner sep=3pt] (L8) at (0, 5.20) {FEX-Emu};
    \node[rectangle, rounded corners=2pt, draw=nodeblue!60, fill=nodeblue!10, minimum width=6cm, minimum height=0.55cm, text=textdark, align=center, inner sep=3pt] (L9) at (0, 5.85) {Legacy32};
    \node[rectangle, rounded corners=2pt, draw=nodelteal!60, fill=nodelteal!10, minimum width=6cm, minimum height=0.55cm, text=textdark, align=center, inner sep=3pt] (L10) at (0, 6.50) {Sysbox};
    \node[rectangle, rounded corners=2pt, draw=nodeorange!60, fill=nodeorange!10, minimum width=6cm, minimum height=0.55cm, text=textdark, align=center, inner sep=3pt] (L11) at (0, 7.15) {WASM};
    \node[rectangle, rounded corners=2pt, draw=nodepurple!60, fill=nodepurple!10, minimum width=6cm, minimum height=0.55cm, text=textdark, align=center, inner sep=3pt] (L12) at (0, 7.80) {Native};

    % Side annotation
    \draw[arrThick] ([xshift=-5mm]L1.west) -- ([xshift=-5mm]L12.north)
      node[midway, left, font=\sffamily\tiny, text=textgray, rotate=90]
        {Security $\rightarrow$};

    % Selection arrow
    \draw[arrBlock, nodeblue] ([xshift=5mm]L5.east) -- ++(1.5,0)
      node[right, font=\sffamily\footnotesize\bfseries, text=nodeblue]
        {Typical on Android};
  \end{tikzpicture}
  \caption{The 12 isolation levels available in Doki, ordered from
    strongest (microVM) to weakest (native). The daemon automatically
    selects the strongest available level.}
  \label{fig:isolation-levels}
\end{figure}
